\begin{theorem}[Cycle operations from an arbitrary actual bond assignment]
\label{thm:peps_actual_cycle_physical_permutation}
\lean{
TNLean.PEPS.exists_unitary_regularActualCyclePhysicalPermutation,
TNLean.PEPS.regularGaugedTreeCycleAssignment_treeGauge_of_internal,
TNLean.PEPS.regularRegionBondExtension_treeGauge,
TNLean.PEPS.regularActualCyclePermutation
}
\leanok
Fix a finite region, an actual spanning tree of its induced graph, and a root.
The bond assignment $u$ determines its normalized tree gauge $k_u$ and its
cycle residual tuple $\omega_u$. Reconstructing those coordinates on internal
edges and retaining the literal exterior recovers $u$ itself.

Let $\phi$ be a permutation of the cycle tuples commuting with simultaneous
conjugation, and suppose the original site maps are $G$-isometric for the
regular virtual action. There is one unitary on the original region spins,
chosen before every $u$ and every boundary configuration, sending the actual
contraction with assignment $u$ to that with internal coefficients reconstructed
from $k_u$ and $\phi(\omega_u)$. Every exterior and crossing operator is unchanged.
The identity extension outside the region acts on the actual closed contraction
in the same way.

This is an auxiliary consequence of the accessible-coordinate and physical
movement constructions of SCP10, local lines 1765--1920 and 2271--2305.
The tree and the conjugation-equivariant permutation remain explicit.
An occupied-to-vacant native plaquette interpretation requires its actual
holonomy calculation; the full four-endpoint braid is not asserted here.
See \cite{gap:rmp_peps_quantum_double_g_isometry}.
\uses{thm:peps_regular_gauged_cycle_permutation}
\end{theorem}
\begin{proof}\leanok
The tree-gradient identity recovers every internal tree coefficient. On a
remaining internal edge, the residual is $k_{\rm head}^{-1}u_e k_{\rm tail}$,
so reconstruction recovers $u_e$ as well. Apply the fixed common-background
physical unitary to these derived coordinates. The literal exterior extension
recovers the input assignment, and the boundary-column and cut identities give
the stated actual contractions.
\end{proof}
